% ==== Bibliography ====
\pdfbookmark[0]{Bibliography}{chap:BibliographyStale}
\yoschapterstar{Bibliography}
\label{chap:BibliographyStale}

\textsc{Achieser}, N. I. 

[$1$] (with I. M. \textsc{Glazman}) Theorie der linearen Operatoren im Hilbert-Raum, Akademie-Verlag $1954$. 

\textsc{Adams}, R. A. [$1$] Sobolev Spaces, Academic Press $1975$. 

Agmon, S. 

[$1$] (with L. \textsc{Nirenberg}) Properties of solutions of ordinary differential equations in Banach space. Comm. P. and Appl. Math. $16$, $121$---$239$ ($1963$). 

\textsc{Aizawa}, S. 

[$1$] A semigroup treatment of the Hamilton-Jacobi equation in one space variable. Hiroshima Math. J. $3$, No. $2$, $367$---$386$ ($1973$). 

\textsc{Akilov}, G. 

[$1$] See \textsc{Kantorovitch}-\textsc{Akilov} [$1$]. 

\textsc{Alexandrov}, P. 

[$1$] (with H. \textsc{Hopf}) Topologie, Vol. I, Springer $1935$. 

\textsc{Aronszajn}, N. 

[$1$] Theory of reproducing kernels. Trans. Amer. Math. Soc. $68$, $337$---$404$ ($1950$). 

\textsc{Balakrishnan}, V. 

[$1$] Fractional powers of closed operators and the semi-groups generated by them. Pacific J. Math. $10$, $419$---$437$ ($1960$). 

\textsc{Banach}, S. 

[$1$] Théorie des Opérations Linéaires, Warszawa $1932$. 

[$2$] Sur la convergence presque partout de fonctionnelles linéaires. Bull. Sci. Math. France $50$, $27$---$32$ and $36$---$43$ ($1926$). 

\textsc{Barbu}, V. 

[$1$] Nonlinear Semigroups and Differential Equations in Banach Spaces, Noordhoff International Publishing $1976$. 

Beboutov, M. 

[$1$] Markoff chains with a compact state space. Rec. Math. $10$, $213$---$238$ ($1942$). 

\textsc{Bénilan}, P. 

[$1$] Solutions intégrales d'équations d'évolution dans un espace de Banach. C. R. Acad. Sci. de Paris $274$, $45$---$50$ ($1972$). 

[$2$] Equations d'évolution dans un espace de Banach quelconque et applications. Thèse, Orsay $1972$. 

\textsc{Bergman}, S. 

[$1$] The Kernel Function and the Conformal Mapping. Mathematical Surveys, No. $5$ ($1950$). 

\textsc{Bers}, L. 

[$1$] Lectures on Elliptic Equations. Summer Seminar in Appl. Math. Univ. of Colorado, $1957$. 

\textsc{Birkhoff}, G. 

[$1$] Lattice Theory. Colloq. Publ. Amer. Math. Soc., $1940$. 

\textsc{Birkhoff}, G. D. 

[$1$] Proof of the ergodic theorem. Proc. Nat. Acad. Sci. USA $17$, $656$---$660$ ($1931$). 

[$2$] (with P. A. \textsc{Smith}) Structure analysis of surface transformations. J. Math. Pures et Appliq. $7$, $345$--$379$ ($1928$). 

Bochner, S. 

[$1$] Integration von Funktionen, deren Wert die Elemente eines Vektorraumes sind. Fund. Math. $20$, $262$--$276$ ($1933$). 

[$2$] Diffusion equations and stochastic processes. Proc. Nat. Acad. Sci. USA $85$, $369$--$370$ ($1949$). 

[$3$] Vorlesungen über Fouriersche Integrale, Akademie-Verlag $1932$. 

[$4$] Beiträge zur Theorie der fastperiodischen Funktionen. Math. Ann. $96$, $119$---$147$ ($1927$). 

Bogolioubov, N. 

[$1$] See \textsc{Krylov}-\textsc{Bogolioubov} [$1$]. 

\textsc{Bohr}, H. 

[$1$] Fastperiodische Funktionen, Springer $1932$. 

\textsc{Bourbaki}, N. 

[$1$] Topologie Générale. Act. Sci. et Ind., nos. $856$, $916$, $1029$, $1045$, $1084$, Hermann $1940$---$42$. 

[$2$] Espaces Vectoriels Topologiques. Act. Sci. et Ind., nos. $1189$, $1229$, Hermann $1953$---$55$. 

\textsc{Brezis}, H. 

[$1$] Opérateurs maximaux monotones et semi-groupes de contractions dans les espaces de Hilbert. Amsterdam: North Holland Publ. Co. $1973$. 

[$2$] (with A. \textsc{Pazy}) Accretive sets and differential equations in Banach spaces. Israel J. of Math. $8$, $367$--$383$ ($1970$). 

\textsc{Browder}, F. E. 

[$1$] Functional analysis and partial differential equations, I---II. Math. Ann. $138$, $55$---$79$ ($1959$) and $145$, $81$---$226$ ($1962$). 

[$2$] Nonlinear Operators and Nonlinear Equations of Evolutions in Banach Spaces, Amer. Math. Soc. $1976$. 

Chacon, R. V. 

[$1$] (with D. S. \textsc{Ornstein}) A general ergodic theorem. Ill. J. Math. $4$, $153$---$160$ ($1960$). 

Choquet, G. 

[$1$] La théorie des représentations intégrales dans les ensembles convexes compacts. Ann. Inst. Fourier$\cdot$$10$, $334$---$344$ ($1960$). 

[$2$] (with P.-A. \textsc{Meyer}) Existence et unicité des représentations intégrales dans les convexes compacts quelconques. Ann. Inst. Fourier $18$, $139$---$154$ ($1963$). 

\textsc{Clarkson}, J. A. 

[$1$] Uniformly convex spaces. Trans. Amer. Math. Soc. $40$, $396$--$414$ ($1936$). \textsc{Crandall}, M. G. 

[$1$] (with A. \textsc{Pazy}) Semi-groups of nonlinear contractions and dissipative sets. J. Funct. Analysis $3$, $376$--$418$ ($1969$). 

[$2$] (with T. \textsc{Liggett}) Generation of semi-groups of nonlinear transformations in general Banach spaces. Amer. J. Math. $93$, $265$--$298$ ($1971$). 

[$3$] The semi-group approach to first order quasilinear equations in several space variables. Israel J. of Math. $12$, $108$--$132$ ($1972$). 

\textsc{Dieudonné}, J. 

[$1$] Recent advances in the theory of locally convex vector spaces. Bull. Amer. math. Soc. $59$, $495$--$512$ ($1953$). 

Doob, J. L. 

[$1$] Stochastic processes with an integral-valued parameter. Trans. Amer. Math. Soc. $44$, $87$--$150$ ($1938$). 

[$2$] Probability theory and the first boundary value problem. Ill. J. Math. $2$, $19$-$36$ ($1958$). 

[$3$] Probability methods applied to the first boundary value problem. Proc. Third Berkeley Symp. on Math. Statist. and Prob. II, $49$---$80$ ($1956$). \textsc{Dorroh}, J. R. 

[$1$] A nonlinear Hille-Yosida-Phillips theorem. J. Funct. Analysis $3$, $345$--$393$ ($1969$). 

\textsc{Dunford}, N. 

[$1$] (with J. \textsc{Schwartz}) Linear Operators, Vol. I, Interscience $1958$. 

[$2$] Uniformity in linear spaces. Trans. Amer. Math. Soc. $44$, $305$-$356$ ($1938$). 

[$3$] On one-parameter groups of linear transformations. Ann. of Math. $39$, $569$-$573$ ($1938$). 

[$4$] (with J. \textsc{Schwartz}) Convergence almost everywhere of operator averages. J. Rat. Mech. Anal. $5$, $129$---$178$ ($1956$). 

[$5$] (with J. \textsc{Schwartz}) Linear Operators, Vol. II, Interscience $1963$. 

[$6$] (with J. \textsc{Schwartz}) Linear Operators, Vol. III, Interscience, to appear. \textsc{Dynkin}, E. B. 

[$1$] Markoff processes and semi-groups of operators. Teorya Veroyatn. $1$ ($1956$). 

[$2$] Infinitesimal operators of Markoff processes. Teorya Veroyatn. $1$ ($1956$). \textsc{Eberlein}, W. F. 

[$1$] Weak compactness in Banach spaces. Proc. Nat. Acad. Sci. USA $88$, $51$--$53$ ($1947$). 

\textsc{Ehrenpreis}, L. 

[$1$] Solutions of some problems of division. Amer. J. Math. $76$, $883$---$903$ ($1954$). 

\textsc{Ekeland}, I. 

[$1$] (with R. \textsc{Temam}) Convex Analysis and Variational Problems, North-Holland Publishing Company $1976$. 

\textsc{Erdélyi}, A. 

[$1$] Operational Calculus and Generalized Functions, Reinhart $1961$. \textsc{Feller}, W. 

[$1$] On the generation of unbounded semi-groups of bounded linear operators. Ann. of Math. $58$, $166$--$174$ ($1953$). 

[$2$] The parabolic differential equation and the associated semi-group of transformations. Ann. of Math. $55$, $468$---$519$ ($1952$). 

[$3$] On the intrinsic form for second order differential operators. Ill. J. Math. $2$, No. $1$, $1$---$18$ ($1958$). 

[$4$] Generalized second order differential operators and their lateral conditions. Ill. J. Math. $1$, No. $4$, $459$--$504$ ($1957$). 

[$5$] Some new connections between probability and classical analysis. Proc. Internat. Congress of Math. $1958$, held at Edinburgh, pp. $69$---$86$. 

[$6$] On differential operators and boundary conditions. Comm. Pure and Appl. Math. $8$, $203$-$216$ ($1955$). 

[$7$] Boundaries induced by non-negative matrices. Trans. Amer. Math. Soc. $88$, $19$---$54$ ($1956$). 

[$8$] On boundaries and lateral conditions for the Kolmogoroff differential equations. Ann. of Math. $65$, $527$-$570$ ($1957$). 

\textsc{Foias}, C. 

[$1$] Remarques sur les semi-groupes distributions d'opérateurs normaux. Portugaliae Math. $19$, $227$---$242$ ($1960$). 

[$2$] See B. Sz. \textsc{Nagy}-C. \textsc{Foias} [$4$]. 

\textsc{Freudenthal}, H. 

[$1$] Über die Friedrichssche Fortsetzung halbbeschränkter Hermitescher Operatoren. Proc. Acad. Amsterdam $89$, $832$--$833$ ($1936$). 

[$2$] Teilweise geordnete Modulen. Proc. Acad. Amsterdam $89$, $641$---$651$ ($1936$). 

\textsc{Friedman}, A. 

[$1$] Generalized Functions and Partial Differential Equations, Prentice-Hall $1963$. 

\textsc{Friedrichs}, K. 

[$1$] Differentiability of solutions of elliptic partial differential equations. Comm. Pure and Appl. Math. $5$, $299$---$326$ ($1953$). 

[$2$] Symmetric positive systems of differential equations. Comm. Pure and Appl. Math. $11$, $333$--$418$ ($1958$). 

[$3$] Spektraltheorie halbbeschränkter Operatoren, I---III. Math. Ann. $109$, $465$---$487$, $685$---$713$ ($1934$) and $110$, $777$---$779$ ($1935$). 

\textsc{Fukamiya}, M. 

[$1$] Topological methods for Tauberian theorem. Tohoku Math. J. $77$---$87$ ($1949$). 

[$2$] See Yosida-Fukamiya [$16$]. 

\textsc{Garding}, L. 

[$1$] Dirichlet's problem for linear elliptic partial differential equations. Math. Scand. $1$, $55$--$72$ ($1953$). 

[$2$] Some trends and problems in linear partial differential equations. Internat. Congress of Math. $1958$, held at Edinburgh, pp. $87$---$102$. 

\textsc{Gelfand}, I. M. 

[$1$] (with G. E. Šilov) Generalized Functions, Vol. I---III, Moscow $1958$, [$2$] Normierte Ringe. Rec. Math. $9$, $3$---$24$ ($1941$). 

[$3$] (with N. Y. \textsc{Vilenkin}) Some Applications of Harmonic Analysis (Vol. IV of Generalized Functions), Moscow $1961$. 

[$4$] (with D. \textsc{Raikov}) On the theory of characters of commutative topological groups. Doklady Akad. Nauk SSSR $28$, $195$---$198$ ($1940$). 

[$5$] (with D. A. \textsc{Raikov} and G. E. \textsc{Šilov}) Commutative Normed Rings, Moscow $1960$. 

\textsc{Gelfand}, I. M. and his collaborators on Generalized Functions: 

[$6$] (with G. E. \textsc{Šilov}) Volume $1$. Properties and Operators, Academic Press $1966$. 

[$7$] (with G. E. \textsc{Šilov}) Volume $2$. Spaces of Fundamental and Generalized Functions, Academic Press $1968$. 

[$8$] (with G. E. \textsc{Šilov}) Volume $3$. Theory of Differential Equations, Academic Press $1967$. 

[$9$] (with N. YA. \textsc{Vilenkin}) Volume $4$. Applications of Harmonic Analysis, Academic Press $1964$. 

[$10$] (with M. I. GRAEV and N. YA. \textsc{Vilenkin}) Volume $5$. Integral Geometry and Representation Theory, Academic Press $1966$. 

\textsc{Hausdorff}, F. 

[$1$] Neue Begründung der Theorie quadratischer Formen von unendlichvielen Veränderlichen. J. reine und angew. Math. $136$, $210$---$271$ ($1909$). 

\textsc{Glazman}, I. M. 

[$1$] See \textsc{Achieser}-\textsc{Glazman} [$1$]. 

\textsc{Grothendieck}, A. 

[$1$] Espaces Vectoriels Topologiques, seconde éd., Sociedade de Mat. de São Paulo $1958$. 

[$2$] Produits Tensoriels Topologiques et Espaces Nucléaires. Memoirs of Amer. Math. Soc. No. $16$ ($1955$). 

\yossub{Hahn, H.}

[$1$] Über Folgen linearer Operatoren. Monatsh. für Math. und Phys. $32$, $3$---$88$ ($1922$). 

[$2$] Über lineare Gleichungssysteme in linearen Räumen. J. reine und angew. Math. $157$, $214$---$229$ ($1927$). 

[$3$] Über die Integrale des Herrn Hellinger und die Orthogonalinvarianten der quadratischen Formen von unendlich vielen Veränderlichen. Monatsh. für Math. und Phys. $28$, $161$---$224$ ($1912$). 

\textsc{Halmos}, P. R. 

[$1$] Measure Theory, van Nostrand $1950$. 

[$2$] Introduction to Hilbert Space and the Theory of Spectral Multiplicity, Chelsea $1951$. 

[$1$] Mengenlehre, W. de Gruyter $1935$. 

\textsc{Hellinger}, E. 

\textsc{Helly}, E. 

[$1$] Über Systeme linearer Gleichungen mit unendlich vielen Unbekannten. Monatsh. für Math. und Phys. $31$, $60$---$91$ ($1921$). 

\yossub{HILBERT, D.}

[$1$] Wesen und Ziele einer Analysis der unendlich vielen unabhängigen Variablen. Rend. Circ. Mat. Palermo $27$, $59$--$74$ ($1909$). 

Hille, E. 

[$1$] (with R. S. \textsc{Phillips}) Functional Analysis and Semi-groups. Colloq. Publ. Amer. Math. Soc., $1957$. It is the second edition of the book below. 

[$2$] Functional Analysis and Semi-groups. Colloq. Publ. Amer. Math. Soc., $1948$. 

[$3$] On the differentiability of semi-groups of operators. Acta Sci. Math. Szeged 12B, $19$---$24$ ($1950$). 

[$4$] On the generation of semi-groups and the theory of conjugate functions. Proc. R. Physiogr. Soc. Lund $21$, $1$---$13$ ($1951$). 

[$5$] Une généralization du problème de Cauchy. Ann. Inst. Fourier $4$, $31$---$48$ ($1952$). 

[$6$] The abstract Cauchy problem and Cauchy's problem for parabolic differential equations. J. d'Analyse Math. $3$, $81$---$196$ ($1954$). 

[$7$] Perturbation methods in the study of Kolmogoroff's equations. Proc. Internat. Congress of Math. $1954$, held at Amsterdam, Vol. III, pp. $365$--$376$. 

[$8$] Linear differential equations in Banach algebras. Proc. Internat. Symposium on Linear Analysis, held at Jerusalem, $1960$, pp. $263$--$273$. 

[$9$] Les probabilités continues en chaine. C. R. Acad. Sci. $230$, $34$---$35$ ($1950$). 

Itô, K. 

\textsc{Hirsch}, F. 

[$1$] Opérateurs codissipatifs. C. R. Acad. Sci. de Paris $270$, $1487$---$1490$ ($1970$). 

[$2$] Familles résolventes générateurs, cogénérateurs, potentiels. Annales L'Institut Fournier de L'Univ. de Grenoble $22$, $89$---$210$ ($1972$). \textsc{Hoffman}, K. 

[$1$] Banach Spaces of Analytic Functions, Prentice-Hall $1962$. \textsc{Hopf}, E. 

[$1$] Ergodentheorie, Springer $1937$. 

[$2$] The general temporally discrete Markoff processes. J. Rat. Mech. and Anal. $8$, $13$---$45$ ($1954$). 

[$3$] On the ergodic theorem for positive linear operators. J. reine und angew. Math. $205$, $101$---$106$ ($1961$). 

\yossub{HOPF, H.}

[$1$] See \textsc{Alexandrov}-\textsc{Hopf} [$1$]. 

\textsc{Hörmander}, L. 

[$1$] On the theory of general partial differential operators. Acta Math. $94$, $161$-$248$ ($1955$). 

[$2$] Lectures on Linear Partial Differential Equations, Stanford Univ. $1960$. 

[$3$] Linear partial differential equations without solutions. Math. Ann. $140$, $169$-$173$ ($1960$). 

[$4$] Local and global properties of fundamental solutions. Math. Scand. $5$, $27$-$39$ ($1957$). 

[$5$] On the interior regularity of the solutions of partial differential equations. Comm. Pure and Appl. Math. $9$, $197$---$218$ ($1958$). 

[$6$] Linear Partial Differential Operators, Springer $1963$. 

\textsc{Hunt}, G. A. 

[$1$] Markoff processes and potentials, I---II. Ill. J. of Math. $1$ and $2$ ($1957$ and $1958$). 

[$1$] (with H. P. McKean) Diffusion Processes and Their Sample Paths, Springer, to appear. 

\yossub{Itô, S.}

[$1$] The fundamental solutions of the parabolic differential equations in differentiable manifold. Osaka Math. J. $5$, $75$--$92$ ($1953$). 

[$2$] (with H. \textsc{Yamabe}) A unique continuation theorem for solutions of a parabolic differential equation. J. Math. Soc. Japan $10$, $314$--$321$ ($1958$). \textsc{Jacobs}, K. 

[$1$] Neuere Methoden und Ergebnisse der Ergodentheorie, Springer $1960$. \textsc{John}, F. 

[$1$] Plane Waves and Spherical Means Applied to Partial Differential Equations, Interscience $1955$. 

\textsc{Kakutani}, S. 

[$1$] Iteration of linear operations in complex Banach spaces. Proc. Imp. Acad. Tokyo $14$, $295$--$300$ ($1938$). 

[$2$] See \textsc{Yosida}-\textsc{Mimura}-\textsc{Kakutani} [$10$]. 

[$3$] Weak topology and regularity of Banach spaces. Proc. Imp. Acad. Tokyo $15$, $169$--$173$ ($1939$). 

[$4$] Concrete representation of abstract $(M)$-spaces. Ann. of Math. $42$, $994$-$1024$ ($1941$). 

[$5$] Concrete representation of abstract $(L)$-spaces and the mean ergodic theorem. Ann. of Math. $42$, $523$-$537$ ($1941$). 

[$6$] Ergodic theorems and the Markoff processes with a stable distribution. Proc. Imp. Acad. Tokyo $16$, $49$--$54$ ($1940$). 

[$7$] See Yosida-Kakutani [$7$]. 

[$8$] Ergodic Theory. Proc. Internat. Congress of Math. $1950$, held at Cambridge, Vol. $2$, pp. $128$--$142$. 

Kantorovitch, L. 

[$1$] (with G. \textsc{Akilov}) Functional Analysis in Normed Spaces, Moscow $1955$. 

Kato, T. 

[$1$] Remarks on pseudo-resolvents and infinitesimal generators of semi-groups. Proc. Japan Acad. $35$, $467$--$468$ ($1959$). 

[$2$] Note on fractional powers of linear operators. Proc. Japan Acad. $36$, $94$---$96$ ($1960$). 

[$3$] Integration of the equation of evolution in a Banach space. J. Math. Soc. of Japan $5$, $208$---$234$ ($1953$). 

[$4$] On linear differential equations in Banach spaces. Comm. Pure and Appl. Math. $9$, $479$-$486$ ($1956$). 

[$5$] Fractional powers of dissipative operators. J. Math. Soc. of Japan $13$, $246$--$274$ ($1961$); II, ibid. $14$, $242$--$248$ ($1962$). 

[$6$] Abstract evolution equations of parabolic type in Banach and Hilbert spaces. Nagoya Math. J. $19$, $93$---$125$ ($1961$). 

[$7$] Fundamental properties of Hamiltonian operators of Schrödinger type. Trans. Amer. Math. Soc. $70$, $195$--$211$ ($1950$). 

[$8$] (with H. \textsc{Tanabe}) On the abstract evolution equation. Osaka Math. J. $14$, $107$---$133$ ($1962$). 

[$9$] Perturbation Theory for Linear Operators, Springer ($1966$). 

[$10$] Semi-groups and temporally inhomogeneous evolution equations, in Equazioni Differenziali Astratte, Centro Internazionale Matematico Estivo, C. I. M. E. $1$ Ciclo, Varenna, $30$ Maggio-$8$ Giugno ($1963$). 

[$11$] Nonlinear semigroups and evolution equations. J. Math. Soc. Japan $19$, $508$-$520$ ($1967$). 

[$12$] Note on the differentiability of nonlinear semigroups. Proc. Symp. Pure Math. Amer. Math. Soc. $16$ ($1970$). 

\textsc{Kelley}, J. L. 

[$1$] General Topology, van Nostrand $1955$. 

[$2$] Note on a theorem of Krein and Milman. J. Osaka Inst. Sci. Tech., Part I, $3$, $1$---$2$ ($1951$). 

\textsc{Kellogg}, O. D. 

[$1$] Foundations of Potential Theory, Springer $1929$. 

\textsc{Khintchine}, A. 

[$1$] Zu Birkhoffs Lösung des Ergodenproblems. Math. Ann. $107$, $485$--$488$ ($1933$). 

Kisyński, J. 

[$1$] Sur les opérateurs de Green des problèmes de Cauchy abstraits. Stud. Math. $23$, $285$---$328$ ($1964$). 

\textsc{Kodaira}, K. 

[$1$] The eigenvalue problem for ordinary differential equations of the second order and Heisenberg's theory of S-matrices. Amer. J. of Math. $71$, $921$--$945$ ($1949$). 

Kolmogorov, A. 

[$1$] Über analytische Methoden in der Wahrscheinlichkeitsrechnung. Math. Ann. $104$, $415$--$458$ ($1931$). 

\textsc{Komatsu}, H. 

[$1$] Abstract analyticity in time and unique continuation property of solutions of a parabolic equation. J. Fac. Sci. Univ. Tokyo, Sect. $1$, $9$, Part $1$, $1$---$11$ ($1961$). 

[$2$] A characterization of real analytic functions. Proc. Japan Acad. $36$, $90$---$93$ ($1960$). 

[$3$] A proof of Kotaké-Narasimhan's theorem. Proc. Japan Acad. $38$, $615$---$618$ ($1962$). 

[$4$] Semi-groups of operators in locally convex spaces. J. Math. Soc. of Japan $16$, $230$--$262$ ($1964$). 

[$5$] Fractional powers of operators, Pacific J. of Math. $19$, No. $2$, $285$--$346$ ($1966$). 

\textsc{Kōmura}, Y. 

[$1$] Nonlinear semigroups in Hilbert space. J. Math. Soc. Japan $19$, $493$---$507$ ($1967$). 

[$2$] Differentiability of nonlinear semigroups. J. Math. Soc. Japan $21$, $375$--$402$ ($1969$). 

\yossub{Konishi, Y.}

[$1$] On the nonlinear semigroups associated with $u_{t} = \Delta \beta(u)$ and $\varphi(u_{t}) = \Delta u$. J. Math. Soc. Japan $25$, $622$-$628$ ($1973$). 

[$2$] On $u_t = u_{xx} - F(u_x)$ and the differentiability of the nonlinear semi-group associated with it. Proc. Japan Acad. $48$, No. $5$, $281$-$286$ ($1972$). \textsc{Kotaké}, T. 

[$1$] Sur l'analyticité de la solution du problème de Cauchy pour certaines classes d'opérateurs paraboliques. C. R. Acad. Sci. Paris $252$, $3716$---$3718$ ($1961$). 

[$2$] (with M. S. \textsc{Narasimhan}) Sur la régularité de certains noyaux associés à un opérateur elliptique. C. R. Acad. Sci. Paris $252$, $1549$---$1550$ ($1961$). \textsc{Köthe}, G. 

[$1$] Topologische lineare Räume, Vol. I, Springer $1960$. 

\yossub{KREIN, M.}

[$1$] (with D. \textsc{Milman}) On extreme points of regularly convex sets. Stud. Math. $9$, $133$--$138$ ($1940$). 

[$2$] (with S. \textsc{Krein}) On an inner characteristic of the set of all continuous functions defined on a bicompact Hausdorff space. Doklady Akad. Nauk SSSR $27$, $429$---$430$ ($1940$). 

\yossub{KREIN, S.}

[$1$] See \textsc{Krein}-\textsc{Krein} [$2$]. 

[$2$] Linear Differential Equations in Banach Spaces, Amer. Math. Soc. $1971$. \textsc{Krylov}, N. 

[$1$] (with N. \textsc{Bogolioubov}) La théorie générale de la mesure dans son application à l'étude des systèmes de la mécanique non linéaires. Ann. of Math. $38$, $65$---$113$ ($1937$). 

\textsc{Ladyzhenskaya}, O. A. 

[$1$] (with I. M. \textsc{Visik}) Problèmes aux limites pour les équations aux dérivées partielles et certaines classes d'équations opérationnelles. Ousp. Mat. Nauk $11$, $41$---$97$ ($1956$). 

\textsc{Lax}, P. D. 

[$1$] (with A. N. \textsc{Milgram}) Parabolic equations, in Contributions to the Theory of Partial Differential Equations, Princeton $1954$. 

\textsc{Maak}, W. 

\textsc{Lewy}, H. 

[$2$] On Cauchy's problem for hyperbolic equations and the differentiability of solutions of elliptic equations. Comm. Pure and Appl. Math. $8$, $615$-$633$ ($1955$). 

[$3$] (with R. S. \textsc{Phillips}) Local boundary conditions for dissipative system of linear partial differential operators. Comm. Pure and Appl. Math. $18$, $427$--$455$ ($1960$). 

[$4$] (with R. S. \textsc{Phillips}) Scattering Theory, Academic Press $1967$. 

\textsc{Leray}, J. 

[$1$] Hyperbolic Differential Equations, Princeton $1952$. 

[$1$] An example of a smooth linear partial differential equation without solutions. Ann. of Math. $66$, $155$-$158$ ($1957$). 

\textsc{Liggett}, T. [$1$] See \textsc{Crandall}-\textsc{Liggett} [$2$]. 

[$1$] Une remarque sur les applications du théorème de Hille-Yosida. J. Math. Soc. Japan $9$, $62$--$70$ ($1957$). 

[$2$] Equations Différentielles Opérationnelles, Springer $1961$. 

[$3$] Espaces d'interpolation et domaines de puissance fractionaires d'opérateurs. J. Math. Soc. Japan $14$, $233$---$241$ ($1962$). 

[$4$] Les semi-groups distributions. Portugaliae Math. $19$, $141$---$164$ ($1960$). 

[$5$] Sur quelques questions d'analyse, de mécanique et de contrôle optimal, Press de Université de Montreal $1976$. 

\textsc{Lumer}, G. 

[$1$] Semi-inner product spaces. Trans. Amer. Math. Soc. $100$, $29$---$43$ ($1961$). 

[$2$] (with R. S. \textsc{Phillips}) Dissipative operators in a Banach space. Pacific J. Math. $11$, $679$--$698$ ($1961$). 

[$1$] Fastperiodische Funktionen. Springer $1950$. 

\textsc{Malgrange}, B. 

[$1$] Existence et approximation des solutions des équations aux dérivées partielles et des équations de convolution. Ann. Inst. Fourier $6$, $271$---$355$ ($1955$---$56$). 

[$2$] Sur une classe d'opérateurs différentielles hypoelliptiques. Bull. Soc. Math. France $58$, $283$---$306$ ($1957$). 

\textsc{Martin}, R. H. 

[$1$] Nonlinear Operators and Differential Equations in Banach Spaces, John Wiley $1976$. 

\textsc{Maruyama}, G. 

[$1$] On strong Markoff property. Mem. Kyushu Univ. $13$ ($1959$). 

\textsc{Masuda}, K. 

[$1$] Hatten Hoteishiki (Evolution Equations), in Japanese, Kinokuniya Book Company $1975$. 

\textsc{Mazur}, S. 

[$1$] Sur les anneaux linéaires. C. R. Acad. Sci. Paris $207$, $1025$---$1027$ ($1936$). 

[$2$] Über konvexe Mengen in linearen normierten Räumen. Stud. Math. $5$, $70$---$84$ ($1933$). 

McKean, H. 

[$1$] See Itô-McKean [$1$]. 

\textsc{Meyer}, P. A. 

[$1$] Probability and Potentials, Blaisdell Publishing Company ($1966$). 

[$2$] See \textsc{Choquet}-\textsc{Meyer} [$2$]. 

Mikusiński, J. [$1$] Operational Calculus, Pergamon $1959$. 

\textsc{Milgram}, A. N. 

[$1$] See \textsc{Lax}-\textsc{Milgram} [$1$]. 

\textsc{Milman}, D. 

[$1$] On some criteria for the regularity of spaces of the type (B). Doklady Akad. Nauk SSSR $20$, $20$ ($1938$). 

[$2$] See \textsc{Krein}-\textsc{Milman} [$1$]. 

\textsc{Mimura}, Y. 

[$1$] See Yosida-Mimura-Kakutani [$10$]. 

[$2$] Über Funktionen von Funktionaloperatoren in einem Hilbertschen Raum. Jap. J. Math $18$, $119$---$128$ ($1936$). 

\textsc{Minlos}, R. A. 

[$1$] Generalized stochastic processes and the extension of measures. Trudy Moscow Math. $8$, $497$--$518$ ($1959$). 

\textsc{Minty}, G. 

[$1$] Monotone (nonlinear) operators in a Hilbert space. Duke Math. J. $29$, $341$-$346$ ($1962$). 

\textsc{Miyadera}, I. 

[$1$] Generation of a strongly continuous semi-group of operators. Tohoku Math. J. $4$, $109$---$121$ ($1952$). 

[$2$] (with S. \textsc{Oharu}) Approximation of semi-groups of nonlinear operators. Tohoku Math. J. $22$, $24$-$47$ ($1970$). 

[$3$] Some remarks on semigroups of nonlinear operators. Tohoku Math. J. $23$, $245$-$258$ ($1971$). 

[$4$] Hisenkei Hangun (Nonlinear Semigroups), in Japanese, Kinokuniya Book Company $1977$. 

Mizohata, S. 

[$1$] Hypoellipticité des équations paraboliques. Bull. Soc. Math. France $85$, $15$--$50$ ($1957$). 

[$2$] Analyticité des solutions élémentaires des systèmes hyperboliques et paraboliques. Mem. Coll. Sci. Univ. Kyoto $32$, $181$---$212$ ($1959$). 

[$3$] Unicité du prolongement des solutions pour quelques opérateurs différentiels paraboliques. Mém. Coll. Sci. Univ. Kyoto, Sér. A $31$, $219$---$239$ ($1958$). 

[$4$] Le problème de Cauchy pour les équations paraboliques. J. Math. Soc. Japan $8$, $269$--$299$ ($1956$). 

[$5$] Systèmes hyperboliques. J. Math. Soc. Japan $11$, $205$---$233$ ($1959$). 

[$6$] The Theory of Partial Differential Equations, Cambridge University Press $1973$. 

\textsc{Morrey}, C. B. 

[$1$] (with L. \textsc{Nirenberg}) On the analyticity of the solutions of linear elliptic systems of partial differential equations. Comm. Pure and Appl. Math. $10$, $271$---$290$ ($1957$). 

\textsc{Nagumo}, M. 

[$1$] Einige analytische Untersuchungen in linearen metrischen Ringen. Jap. J. Math. $18$, $61$---$80$ ($1936$). 

[$2$] \operatorname{Re}-topologization of functional spaces in order that a set of operators will be continuous. Proc. Japan Acad. $37$, $550$---$552$ ($1961$). 

\textsc{Nagy}, B. von Sz. 

[$1$] Spektraldarstellung linearer Transformationen des Hilbertschen Raumes, Springer $1942$. 

[$2$] See Riesz-\textsc{Nagy} [$3$]. 

[$3$] Prolongements des transformations de l'espace de Hilbert qui sortent de cet espace. Akad. Kiado, Budapest $1955$. 

[$4$] (with C. \textsc{Foias}) Harmonic Analysis of Operators on Hilbert spaces, North-Holland Publishing Company $1970$. 

Naimark, M. A. 

[$1$] Normed Rings, P. Noordhoff $1959$. 

[$2$] Lineare differentiale Operatoren, Akad. Verlag $1960$. 

[$3$] Über Spektralfunktionen eines symmetrischen Operators. Izvestia Akad. Nauk SSSR $17$, $285$--$296$ ($1943$). 

Nakano, H. 

[$1$] Unitärinvariante hypermaximale normale Operatoren, Ann. of Math. $42$, $657$--$664$ ($1941$). 

\textsc{Narasimhan}, M. S. 

[$1$] See \textsc{Kotaké}-\textsc{Narasimhan} [$2$]. 

\textsc{Nelson}, E. 

[$1$] Analytic vectors. Ann. of Math. $670$, $572$---$615$ ($1959$). 

[$2$] Feynman integrals and the Schrödinger equations. J. of Math. Physics $5$, $332$--$343$ ($1964$). 

\textsc{Neumann}, J. \textsc{von} 

[$1$] Allgemeine Eigenwerttheorie Hermitescher Funktionaloperatoren. Math. Ann. $102$, $49$---$131$ ($1929$). 

[$2$] On rings of operators, III. Ann. of Math. $41$, $94$---$161$ ($1940$). 

[$3$] Zur Operatorenmethode in der klassischen Mechanik. Ann. of Math. $33$, $587$--$643$ ($1932$). 

[$4$] Almost periodic functions in a group. I. Trans. Amer. Math. Soc. $36$, $445$-$492$ ($1934$). 

[$5$] Über adjungierte Funktionaloperatoren. Ann. of Math. $33$, $249$---$310$ ($1932$). 

[$6$] Über die analytischen Eigenschaften von Gruppen linearer Transformationen und ihrer Darstellungen. Math. Z. $30$, $3$---$42$ ($1929$). 

[$7$] Zur Algebra der Funktionaloperatoren und Theorie der normalen Operatoren. Math. Ann. $102$, $370$--$427$ ($1929$--$30$). 

[$8$] Über einen Satz von Herrn M. H. \textsc{Stone}. Ann. of Math. $33$, $567$---$573$ ($1932$). 

\textsc{Nirenberg}, L. 

[$1$] Remarks on strongly elliptic partial differential equations. Comm. Pure and Appl. Math. $8$, $643$--$674$ ($1955$). 

[$2$] On elliptic partial differential equations. Ann. Scuola Norm. Sup. Pisa $18$, $115$--$162$ ($1959$). 

[$3$] See \textsc{Morrey}-\textsc{Nirenberg} [$1$]. 

[$4$] See \textsc{Agmon}-\textsc{Nirenberg} [$1$]. 

\textsc{Oharu}, S. 

[$1$] See \textsc{Miyadera}-\textsc{Oharu} [$2$]. 

\textsc{Okamoto}, S. See \textsc{Yosida}-\textsc{Okamoto} [$40$]. 

\textsc{Ornstein}, D. S. 

[$1$] See \textsc{Chacon}-\textsc{Ornstein} [$1$]. 

\textsc{Paley}, R. E. A. C. 

[$1$] (with N. \textsc{Wiener}) Fourier Transforms in the Complex Domain, Colloq. Publ. Amer. Math. Soc., $1934$. 

Pazy, A. 

[$1$] See \textsc{Brezis}-\textsc{Pazy} [$2$]. 

[$2$] See \textsc{Crandall}-\textsc{Pazy} [$1$]. 

\textsc{Peetre}, J. 

[$1$] A proof of the hypoellipticity of formally hypoelliptic differential operators. Comm. Pure and Appl. Math. $16$, $737$--$747$ ($1961$). 

\textsc{Peter}, F. 

[$1$] (with H. \textsc{Weyl}) Die Vollständigkeit der primitiven Darstellungen einer geschlossenen kontinuierlichen Gruppe. Math. Ann. $97$, $737$--$755$ ($1927$). 

\textsc{Petrowsky}, I. G. 

[$1$] Sur l'analyticité des solutions d'équations différentielles. Rec. Math. $47$ ($1939$). 

Pettis, B. J. 

[$1$] On integration in vector spaces. Trans. Amer. Math. Soc. $44$, $277$---$304$ ($1938$). 

\textsc{Phelps}, R. R. 

[$1$] Lectures on Choquet's Theorem, van Nostrand $1966$. 

\textsc{Phillips}, R. S. 

[$1$] See \textsc{Hille}-\textsc{Phillips} [$1$]. 

[$2$] The adjoint semi-group. Pacific J. Math. $5$, $269$---$283$ ($1955$). 

[$3$] An inversion formula for Laplace transform and semi-groups of linear operators. Ann. of Math. $59$, $325$-$356$ ($1954$). 

[$4$] See \textsc{Lumer}-\textsc{Phillips} [$2$]. 

[$5$] On the generation of semi-groups of linear operators. Pacific J. Math. $2$, $343$-$369$ ($1952$). 

[$6$] On the integration of the diffusion equation with boundaries. Trans. Amer. Math. Soc. $98$, $62$--$84$ ($1961$). 

[$7$] Dissipative operators and parabolic partial differential operators. Comm. Pure and Appl. Math. $12$, $249$---$276$ ($1959$). 

[$8$] Dissipative hyperbolic systems. Trans. Amer. Math. Soc. $86$, $109$---$173$ ($1957$). 

[$9$] Dissipative operators and hyperbolic systems of partial differential equations. Trans. Amer. Math. Soc. $90$, $193$--$254$ ($1959$). 

[$10$] See \textsc{Lax}-\textsc{Phillips} [$3$]. 

[$11$] See P. D. \textsc{Lax}-R. S. \textsc{Phillips} [$4$]. 

Pitt, H. R. 

[$1$] Tauberian Theorems. Tata Inst. of Fund. Research, $1958$. 

\textsc{Pontrjagin}, L. 

[$1$] Topological Groups, Princeton $1939$. 

\textsc{Poulsen}, E. T. 

[$1$] Evolutionsgleichungen in Banach-Räumen, Math. Zeitschr. $90$, $286$--$309$ ($1965$). 

Quinn, B. K. 

[$1$] Solution with shocks: An example of an $\mathbf{L}_1$-contractive semigroup. Comm. Pure and Appl. Math. $24$, $125$-$132$ ($1971$). 

\textsc{Rado}, T. 

[$1$] Subharmonic Functions, Springer $1937$. 

\textsc{Raikov}, D. A. 

[$1$] See \textsc{Gelfand}-\textsc{Raikov} [$4$]. 

[$2$] See \textsc{Gelfand}-\textsc{Raikov}-\textsc{Šilov} [$5$]. 

Riesz, F. 

\textsc{Rasmussen}, S. 

[$1$] Non-linear semi-groups, evolution equations and product integral representations. Various Publication Series, No. $2$, Aarhus Universitet ($1971$/$72$). 

Ray, D. 

[$1$] Resolvents, transition functions and strongly Markovian processes. Ann. of Math. $70$, $43$--$72$ ($1959$). 

Rickart, C. E. 

[$1$] General Theory of Banach Algebras, van Nostrand $1960$. 

[$1$] Zur Theorie des Hilbertschen Raumes. Acta Sci. Math. Szeged $7$, $34$---$38$ ($1934$). 

[$2$] Über lineare Funktionalgleichungen. Acta Math. $41$, $71$---$98$ ($1918$). 

[$3$] (with B. von Sz. \textsc{Nagy}) Leçons d'Analyse Fonctionelle, Akad. Kiado, Budapest $1952$. 

[$4$] Some mean ergodic theorems. J. London Math. Soc. $13$, $274$---$278$ ($1938$). 

[$5$] Sur les fonctions des transformations hermitiennes dans l'espace de Hilbert. Acta Sci. Math. Szeged $7$, $147$---$159$ ($1935$). 

[$6$] Sur la Décomposition des Opérations Linéaires. Proc. Internat. Congress of Math. $1928$, held at Bologna, Vol. III, $143$--$148$. 

\textsc{Rockafellar}, R. T. 

[$1$] Convex Analysis, Princeton University Press $1972$. 

Ryll-Nardzewski, C. 

[$1$] See J. Mikusiński [$1$]. 

\yossub{Saks, S.}

[$1$] Theory of the Integral, Warszawa $1937$. 

[$2$] Addition to the note on some functionals. Trans. Amer. Math. Soc. $85$, $967$-$974$ ($1933$). 

Sato, K. 

[$1$] Positive pseudo-resolvents in Banach lattices. J. Fac. Science, The Univ. of Tokyo, Sect. IA, Vol. $17$, Nos. $1$---$2$, $305$---$313$ ($1970$). 

[$2$] Potential operators for Markov processes, to appear. 

Sato, M. 

[$1$] Theory of hyperfunctions I, J. Fac. Sci. Univ. Tokyo, Sect. $1$, $8$, $139$--$193$ ($1959$); II, ibid. $8$, $387$--$437$ ($1960$). 

[$2$] (with T. KAWAI and M. KASHIWARA) Hyperfunctions and Pseudo-Differential Equations, Lecture Notes in Mathematics $287$, Springer $1973$. 

[$3$] (with many authors) Recent Development in Hyperfunction Theory and its Application to Physics (Microlocal Analysis of S-Matrices and Related Quantities), Lecture Notes in Physics $39$, Springer $1975$. 

\textsc{Schapira}, P. 

[$1$] Théorie des Hyperfonctions, Lecture Notes in Mathematics $126$, Springer $1970$. 

\textsc{Schatten}, R. 

[$1$] A Theory of Cross-spaces, Princeton $1950$. 

Schauder, J. 

[$1$] Über lineare, vollstetige Funktionaloperationen. Stud. Math. $2$, $1$---$6$ ($1930$). 

Schwartz, J. 

[$1$] See \textsc{Dunford}-\textsc{Schwartz} [$1$]. 

[$2$] See \textsc{Dunford}-\textsc{Schwartz} [$4$]. 

[$3$] See \textsc{Dunford}-\textsc{Schwartz} [$5$]. 

[$4$] See \textsc{Dunford}-\textsc{Schwartz} [$6$]. 

Schwartz, L. 

[$1$] Théorie des Distributions, vol. I et II, Hermann $1950$, $1951$. 

[$2$] Transformation de Laplace des distributions. Comm. Sém. Math. de l'Univ. de Lund, tome suppl. dédié à M. \textsc{Riesz}, $196$---$206$ ($1952$). 

[$3$] Lectures on Mixed Problems in Partial Differential Equations and the Representation of Semi-groups. Tata Inst. Fund. Research, $1958$. 

[$4$] Les équations d'évolution liées au produit de compositions. Ann. Inst. Fourier $2$, $165$---$169$ ($1950$---$1951$). 

[$5$] Exposé sur les travaux de Gårding, Séminaire Bourbaki, May $1952$. 

[$6$] Théorie des Distributions, Hermann $1966$. 

[$7$] Radon Measures on Arbitray Topological Spaces and Cylindrical Measures, Tata Institute of Fundamental Research and Oxford University Press $1973$. 

SEgal, I. E. 

[$1$] The span of the translations of a function in a Lebesgue space. Proc. Nat. Acad. Sci. USA $80$, $165$--$169$ ($1944$). 

\textsc{Shmulyan}, V. L. 

[$1$] Über lineare topologische Räume. Math. Sbornik, N. S. $7$ ($49$), $425$--$448$ ($1940$). 

Šilov, G. 

[$1$] See \textsc{Gelfand}-\textsc{Šilov} [$1$]. 

[$2$] See GENFAD-RALKOV-\textsc{Šilov} [$5$]. 

\textsc{Smith}, P. A. 

[$1$] See \textsc{Birkhoff}-\textsc{Smith} [$2$]. 

Sobolev, S. L. 

[$1$] Sur un théorème d'analyse fonctionnelle. Math. Sbornik $45$, $471$---$496$ ($1938$). 

[$2$] Certaines Applications de l'Analyse Fonctionnelle à la Physique Mathématique, Leningrad $1945$. 

Sobolevski, P. E. 

[$1$] Parabolic type equations in Banach spaces. Trudy Moscow Math. $10$, $297$---$350$ ($1961$). 

\textsc{Stein}, E. M. 

[$1$] (with G. \textsc{Weiss}) Introduction to Fourier Analysis on Euclidean Spaces, Princeton University Press $1971$. 

[$2$] Singular Integrals and Differentiability Properties of Functions, Princeton University Press $1970$. 

\textsc{Stone}, M. H. 

[$1$] Linear Transformations in Hilbert Space and Their Applications to Analysis. Colloq. Publ. Amer. Math. Soc., $1932$. 

[$2$] On one-parameter unitary groups in Hilbert space. Ann. of Math. $33$, $643$-$648$ ($1932$). 

Szegö, G. 

[$1$] Orthogonal Polynomials. Colloq. Pub. Amer. Math. Soc., $1948$. 

\textsc{Tanabe}, H. 

[$1$] Evolution equations of parabolic type. Proc. Japan Acad. $37$, $610$---$613$ ($1961$). 

[$2$] On the equations of evolution in a Banach space. Osaka Math. J. $12$, $365$-$613$ ($1960$). 

[$3$] A class of the equations of evolution in a Banach space. Osaka Math. J. $11$, $121$---$145$ ($1959$). 

[$4$] Remarks on the equations of evolution in a Banach space. Osaka Math. J. $12$, $145$--$166$ ($1960$). 

[$5$] See \textsc{Kato}-\textsc{Tanabe} [$8$]. 

[$6$] Hatten Hoteishiki (Evolution Equations), in Japanese, Iwanami Shoten $1975$. 

\textsc{Tannaka}, T. 

[$1$] Dualität der nicht-kommutativen bikompakten Gruppen. Tohoku Math. J. $53$, $1$---$12$ ($1938$). 

\textsc{Taylor}, A. 

[$1$] Introduction to Functional Analysis, Wiley $1958$. 

\textsc{Temam}, R. 

[$1$] See I. \textsc{Ekeland}-R. \textsc{Temam} [$1$]. 

\textsc{Titchmarsh}, E. C. 

[$1$] Introduction to the Theory of Fourier Integrals, Oxford $1937$. 

[$2$] Eigenfunction Expansion Associated with Second-order Differential Equations, Vol. I---II, Oxford $1946$---$1958$. 

Trèves, F. 

[$1$] Lectures on Partial Differential Equations with Constant Coefficients. Notas de Matematica, No. $7$, Rio de Janeiro $1961$. 

\textsc{Trotter}, H. F. 

[$2$] Basic Linear Differential Equations, Academic Press $1975$. 

[$1$] Approximation of semi-groups of operators. Pacific J. Math. $8$, $887$---$919$ ($1958$). 

\textsc{Vilenkin}, N. Y. 

[$1$] See \textsc{Gelfand}-\textsc{Vilenkin} [$3$]. 

Visik, I. M. 

[$1$] See \textsc{Ladyzhenskaya}-\textsc{Visik} [$1$]. 

\textsc{Vitali}, G. 

[$1$] Sull'integrazioni per serie. Rend. Circ. Mat. di Palermo $23$, $137$---$155$ ($1907$). 

\textsc{Watanabe}, J. 

[$1$] On some properties of fractional powers of linear operators. Proc. Japan Acad. $87$, $273$---$275$ ($1961$). 

\textsc{Webb}, G. 

[$1$] Representation of semigroups of nonlinear nonexpansive transformations in Banach spaces. J. Math. Mech. $19$, $157$--$170$ ($1969$). 

\textsc{Wecken}, F. J. 

[$1$] Unitärinvariante selbstadjungierte Operatoren. Math. Ann. $116$, $422$---$455$ ($1939$). 

\textsc{Weil}, A. 

[$1$] Sur les fonctions presque périodiques de von Neumann. C. R. Acad. Sci. Paris $200$, $38$--$40$ ($1935$). 

\textsc{Weiss}, G. 

[$1$] See E. M. \textsc{Stein}-G. \textsc{Weiss} [$1$]. 

\textsc{Weyl}, H. 

[$1$] The method of orthogonal projection in potential theory. Duke Math. J. $7$, $414$---$444$ ($1940$). 

[$2$] Über gewöhnliche Differentialgleichungen mit Singularitäten und die zugehörigen Entwicklungen willkürlicher Funktionen. Math. Ann. $68$, $220$---$269$ ($1910$). 

[$3$] See \textsc{Peter}-\textsc{Weyl} [$1$]. 

Wiener, N. 

[$1$] See \textsc{Paley}-\textsc{Wiener} [$1$]. 

[$2$] Tauberian theorems. Ann. of Math. $33$, $1$---$100$ ($1932$). 

[$3$] The Fourier Integral and Certain of Its Applications, Cambridge $1933$. 

\yossub{YAMABE, H.}

[$1$] See Ito-Yamabe [$2$]. 

\textsc{Yamada}, A. 

[$1$] On the correspondence between potential operators and semi-groups associated with Markov processes. Zeitschr. für Wahrscheinlichkeitstheorie und verw. Geb. $15$, $230$---$238$ ($1970$). 

Yosida, K. 

[$1$] Lectures on Differential and Integral Equations, Interscience $1960$. 

[$2$] Vector lattices and additive set functions. Proc. Imp. Acad. Tokyo $17$, $228$-$232$ ($1940$). 

[$3$] Mean ergodic theorem in Banach spaces. Proc. Imp. Acad. Tokyo $14$, $292$-$294$ ($1938$). 

[$4$] Ergodic theorems for pseudo-resolvents. Proc. Japan Acad. $37$, $422$---$425$ ($1961$). 

[$5$] On the differentiability and the representation of one-parameter semigroups of linear operators. J. Math. Soc. Japan $1$, $15$--$21$ ($1948$). 

[$6$] Holomorphic semi-groups in a locally convex linear topological space. Osaka Math. J., $15$, $51$--$57$ ($1963$). 

[$7$] (with S. \textsc{Kakutani}) Operator-theoretical treatment of Markoff process and mean ergodic theorems. Ann. of Math. $42$, $188$--$228$ ($1941$). 

[$8$] Fractional powers of infinitesimal generators and the analyticity of the semi-groups generated by them. Proc. Japan Acad. $36$, $86$--$89$ ($1960$). 

[$9$] Quasi-completely continuous linear functional operators. Jap. J. Math. $15$, $297$---$301$ ($1939$). 

[$10$] (with Y. \textsc{Mimura} and S. \textsc{Kakutani}) Integral operators with bounded measurable kernel. Proc. Imp. Acad. Tokyo $14$, $359$---$362$ ($1938$). 

[$11$] On the group embedded in the metrical complete ring. Jap. J. Math. $13$, $7$---$26$ ($1936$). 

[$12$] Normed rings and spectral theorems. Proc. Imp. Acad. Tokyo $19$, $356$---$359$ ($1943$). 

[$13$] On the unitary equivalence in general euclid spaces. Proc. Jap. Acad. $22$, $242$--$245$ ($1946$). 

[$14$] On the duality theorem of non-commutative compact groups. Proc. Imp. Acad. Tokyo $19$, $181$---$183$ ($1943$). 

[$15$] An abstract treatment of the individual ergodic theorems. Proc. Imp. Acad. Tokyo $16$, $280$---$284$ ($1940$). 

[$16$] (with M. \textsc{Fukamiya}) On vector lattice with a unit, II. Proc. Imp. Acad. Tokyo $18$, $479$---$482$ ($1941$). 

[$17$] Markoff process with a stable distribution. Proc. Imp. Acad. Tokyo $16$, $43$--$48$ ($1940$). 

[$18$] Ergodic theorems of Birkhoff-Khintchine's type. Jap. J. Math. $17$, $31$-$36$ ($1940$). 

[$19$] Simple Markoff process with a locally compact phase space. Math. Japonicae $1$, $99$---$103$ ($1948$). 

[$20$] Brownian motion in a homogeneous Riemannian space. Pacific J. Math. $2$, $263$---$270$ ($1952$). 

[$21$] An abstract analyticity in time for solutions of a diffusion equation. Proc. Japan Acad. $35$, $109$--$113$ ($1959$). 

[$22$] An operator-theoretical integration of the wave equation. J. Math. Soc. Japan $8$, $79$---$92$ ($1956$). 

[$23$] Semi-group theory and the integration problem of diffusion equations. Internat. Congress of Math. $1954$, held at Amsterdam, Vol. $3$, pp. $864$---$873$. 

[$24$] On the integration of diffusion equations in Riemannian spaces. Proc. Amer. Math. Soc. $3$, $864$-$873$ ($1952$). 

[$25$] On the fundamental solution of the parabolic equations in a Riemannian space. Proc. Amer. Math. Soc. $3$, $864$--$873$ ($1952$). 

[$26$] An extension of Fokker-Planck's equation. Proc. Japan Acad. $25$, $1$---$3$ ($1949$). 

[$27$] Brownian motion on the surface of $3$-sphere. Ann. of Math. Statist. $20$, $292$-$296$ ($1949$). 

[$28$] On the integration of the equation of evolution. J. Fac. Sci. Univ. Tokyo, Sect. $1$, $9$, Part $5$, $397$---$402$ ($1963$). 

[$29$] An operator-theoretical treatment of temporally homogeneous Markoff processes. J. Math. Soc. Japan $1$, $244$---$253$ ($1949$). 

[$30$] On holomorphic Markov processes, Proc. Japan Acad. $42$, No. $4$, $313$--$317$ ($1966$). 

[$31$] A perturbation theorem for semi-groups of linear operators, Proc. Japan Acad. $41$, No. $8$, $645$--$647$ ($1965$). 

[$32$] Postitive resolvents and potentials (An operator-theoretical treatment of Hunt's theory of potentials), Zeitschr. für Wahrscheinlichkeitstheorie und verw. Geb. $8$, $210$--$218$ ($1967$). 

[$33$] Time dependent evolution equations in a locally convex space, Math. Zeitschr. $162$, $83$-$86$ ($1965$). 

[$34$] The existence of the potential operator associated with an equi-continuous semi-group of class $(C_0)$. Studia Math. $31$, $531$-$533$ ($1968$). 

[$35$] On the existence and a characterization of abstract potential operators. Proc. Colloq. Functional Analysis, Liège ($1970$), $129$--$136$. 

[$36$] (with T. \textsc{Watanabe} and H. \textsc{Tanaka}) on the pre-closedness of the potential operators. J. Math. Soc. Japan $20$, $419$---$421$ ($1968$). 

[$37$] The pre-closedness of Hunt's potential operators and its applications. Proc. Intern. Conf. Functional Analysis and Related Topics, Toyko, $324$---$331$ ($1969$). 

[$38$] On the potential operators associated with Brownian motions. J. d'Analyse Mathématique $23$, $461$---$465$ ($1970$). 

[$39$] Abstract potential operators on Hilbert space. Publ. of the Research Inst. for Math. Sci., Kyoto Univ. $8$, No. $1$, $201$---$205$ ($1972$). 

[$40$] (with S. \textsc{Okamoto}) A note on Mikusiński's operational calculus. Proc. Japan Acad. $56$, Ser. A, $1$---$3$ ($1980$). 

\textsc{Yushkevitch}, A. A. 

[$1$] On strong Markoff processes. Teorya Veroyatn. $2$ ($1957$). 
